\section{从镜像最优性方程到四种可训练目标}
\label{sec:ch2-objective-derivation}

定理 \ref{thm:ch2-gibbs} 给出的不是孤立的分布结论，而是一条可以被 inner
optimizer 直接拟合的逐 token 方程。下面先把这条方程、平方损失和 BPO 的
中心化残差逐步写出，再进入四种算法的实现。

\subsection{镜像最优性方程与训练测度}
\label{subsec:ch2-optimality-bridge}

把式 \eqref{eq:ch2-optimal-policy} 移项，得到
\begin{equation}
\label{eq:ch2-logratio-equation}
\log\frac{ p_k^{\mathrm G}(a\mid s)}{ q_k(a\mid s)}
=\eta f(s,a)-\log Z_f(s).
\end{equation}
在参数化 inner 中，左侧由 $ p_{k,n}$ 近似。因此定义镜像方程残差
\begin{equation}
R_{k,n}^{\mathrm{Part}}(s,a)
\coloneqq
\log\frac{ p_{k,n}(a\mid s)}{ q_k(a\mid s)}
-\eta f(s,a)+\log Z_f(s).
\label{eq:ch2-mirror-equation-residual}
\end{equation}
精确 mirror 解的条件是 $R_{k,n}^{\mathrm{Part}}(s,a)=0$ 对所有有正概率的
$(s,a)$ 成立。

本书选择 $L_2(d_kq_k)$ 作为平方回归的投影几何，以便和冻结 rollout 的采样测度对齐。原始 mirror objective 使用 student 测度；冻结 anchor 本身并不强迫平方回归只能采用这一范数。其他严格正权重在完整可实现条件下也能恢复同一零残差分布，表达受限时的投影与梯度则会改变。所选范数定义为
\begin{equation}
\lVert R_{k,n}^{\mathrm{Part}}\rVert_k^2
\coloneqq
\mathbb E_{s\sim d_k}\mathbb E_{a\sim q_k(\cdot\mid s)}
\left[R_{k,n}^{\mathrm{Part}}(s,a)^2\right].
\label{eq:ch2-anchor-projection-norm}
\end{equation}
这也是 on-policy rollout 下可直接估计的投影几何。若改用行为测度 $\beta_k$
采样，必须乘以冻结的 importance weight
$ q_k(a\mid s)/\beta_k(a\mid s)$，否则求解的是 $L_2(d_k\beta_k)$ 投影。

\subsection{PMD-Part 与 PMD-Mean：从方程到平方损失}
\label{subsec:ch2-pmd-loss-derivation}

PMD-Part 直接最小化式 \eqref{eq:ch2-mirror-equation-residual} 的平方。将
student log-ratio 和冻结 target 展开，得到
\begin{equation}
\label{eq:ch2-regression-from-mirror}
\mathcal L_{\mathrm{Part}}(\theta)
=\frac12\mathbb E_{s\sim d_k,a\sim q_k}
\left[
\log\frac{ p_{k,n}(a\mid s)}{ q_k(a\mid s)}
-\eta f(s,a)+\operatorname{sg}\!\left(\log Z_f(s)\right)
\right]^2.
\end{equation}
其中 $\operatorname{sg}$ 表示停止梯度；$ q_k$、$f$ 和 $Z_f$ 在本轮
inner 中都冻结。因为
\begin{equation*}
\nabla_\theta\log\frac{ p_{k,n}(a\mid s)}{ q_k(a\mid s)}
=\nabla_\theta\log p_{k,n}(a\mid s),
\end{equation*}
其梯度为
\begin{equation}
\nabla_\theta\mathcal L_{\mathrm{Part}}
=\mathbb E_{d_k, q_k}
\left[\operatorname{sg}\!\left(
\log\frac{ p_{k,n}}{ q_k}-\eta f+\log Z_f
\right)\nabla_\theta\log p_{k,n}(a\mid s)\right].
\label{eq:ch2-part-gradient-bridge}
\end{equation}
这说明 PMD-Part 是 anchor 测度下对完整 mirror log-ratio 的 $L_2$ 投影，
而不是把一个 teacher label 当作普通的 supervised token label。

PMD-Mean 使用原论文的均值近似，而不是对 student log-ratio 做条件中心化。使用第一章定义的冻结 target $G_k(s,a)=\eta(f(s,a)-\bar f_k(s))$。
将完整 target 中的 $\log Z_f(s)$ 用 $\eta\bar f_k(s)$ 近似，得到
\begin{equation}
\mathcal L_{\mathrm{Mean}}(\theta)
=\frac12\mathbb E_{d_k, q_k}
\left[\log\frac{ p_{k,n}(a\mid s)}{ q_k(a\mid s)}
-\eta\bigl(f(s,a)-\bar f_k(s)\bigr)\right]^2.
\label{eq:ch2-mean-loss-bridge}
\end{equation}
这里的 reward center 在整个 inner loop 冻结；损失中没有
$\operatorname{KL}( q_k\Vert p_{k,n})$。定义 $R_{k,n}^{\mathrm{Mean}}=U_{k,n}-G_k$，直接求导得到
\begin{equation}
\nabla_\theta\mathcal L_{\mathrm{Mean}}
=\mathbb E_{d_k, q_k}
\left[(U_{k,n}-G_k)\nabla_\theta\log p_{k,n}(a\mid s)\right].
\label{eq:ch2-mean-gradient-bridge}
\end{equation}
均值近似遗漏的归一化修正是
\begin{equation*}
\log Z_f(s)-\eta\bar f_k(s)
=\log\mathbb E_{ q_k}[e^{G_k(s,\cdot)}\mid s]\ge0.
\end{equation*}
非恒定信号对应严格正值，此时 $G_k$ 不满足指数归一化约束。
因此 PMD-Mean 即使具有完整函数类和无限样本，也不恢复原始 Gibbs 解；其受约束总体解见定理~\ref{thm:ch2-mean-lambert}。
若改为拟合 $U_{k,n}-\mathbb E_{ q_k}[U_{k,n}\mid s]-G_k$，得到的是后面的 BPO 中心化残差，已经改变了 PMD-Mean 的目标。

\subsection{BPO：中心化最优性残差}
\label{subsec:ch2-bpo-loss-derivation}

BPO 将镜像方程残差按冻结 anchor 条件中心化。配分函数项是 action-independent 常数，中心化后消失：
\begin{align}
R_{k,n}^{\mathrm{BPO}}(s,a)
&\coloneqq U_{k,n}(s,a)-\eta f(s,a)
-\mathbb E_{ q_k}[U_{k,n}(s,\cdot)-\eta f(s,\cdot)\mid s]
\label{eq:ch2-bpo-from-optimality}\\
&=U_{k,n}(s,a)-G_k(s,a)+\mathcal K_{k,n}^{S}(s).
\label{eq:ch2-bpo-expanded}
\end{align}
这里 $\mathbb E_{ q_k}[U_{k,n}\mid s]=-\mathcal K_{k,n}^{S}$，所以 student KL 前是正号。OPD 中 $f=r_k$，故
\begin{equation}
R_{k,n}^{\mathrm{BPO}}(s,a)
=\log\frac{ p_{k,n}(a\mid s)}{ q_k(a\mid s)}+\mathcal K_{k,n}^{S}(s)
-\eta\left(\log\frac{ t(a\mid s)}{ q_k(a\mid s)}+\mathcal K_k^{T}(s)\right).
\label{eq:ch2-bpo-opd-expanded}
\end{equation}
BPO 最小化
\begin{equation}
\mathcal L_{\mathrm{BPO}}(\theta)
=\frac1{2\eta}\mathbb E_{d_k, q_k}[(R_{k,n}^{\mathrm{BPO}})^2].
\label{eq:ch2-bpo-loss-bridge}
\end{equation}
使用第一章定义的 score $\psi_{k,n}$ 及其 anchor 均值 $\bar\psi_{k,n}$，保留 student KL 的导数得到
\begin{equation*}
\nabla_\theta\mathcal K_{k,n}^{S}=-\bar\psi_{k,n},\qquad
\nabla_\theta R_{k,n}^{\mathrm{BPO}}=\psi_{k,n}-\bar\psi_{k,n}.
\end{equation*}
因此精确梯度为
\begin{align}
\nabla_\theta\mathcal L_{\mathrm{BPO}}
&=\frac1\eta\mathbb E_{d_k, q_k}[R_{k,n}^{\mathrm{BPO}}(\psi_{k,n}-\bar\psi_{k,n})]
\label{eq:ch2-bpo-gradient-bridge}\\
&=\mathbb E_{d_k, q_k}\left[
\left(\frac{U_{k,n}+\mathcal K_{k,n}^{S}}\eta-f+\bar f_k\right)\psi_{k,n}\right].
\end{align}
第二行使用 $\mathbb E_{ q_k}[R_{k,n}^{\mathrm{BPO}}\mid s]=0$，已把 $\eta^{-1}$ 乘入权重。梯度下降使用上述梯度的负值。消失的是 residual 均值乘以 score 均值，score 的 anchor 期望 $\bar\psi_{k,n}$ 本身通常不为零。

\subsection{MDPO：最小化原始 mirror objective 的负值}
\label{subsec:ch2-mdpo-loss-derivation}

MDPO 使用原 mirror objective 的负值作为损失，统一执行梯度下降：
\begin{equation}
\mathcal L_{\mathrm{MDPO}}(\theta)
=-\mathbb E_{s\sim d_k,a\sim p_{k,n}}[f(s,a)]
+\frac1\eta\mathbb E_{s\sim d_k}\operatorname{KL}( p_{k,n}(\cdot\mid s)\Vert q_k(\cdot\mid s)).
\label{eq:ch2-mdpo-objective-bridge}
\end{equation}
逐历史应用式~\eqref{eq:ch2-kl-decomp}，得到
\begin{equation}
\mathcal L_{\mathrm{MDPO}}(\theta\mid s)
=-\frac1\eta\log Z_f(s)
+\frac1\eta\operatorname{KL}( p_{k,n}(\cdot\mid s)\Vert p_k^{\mathrm G}(\cdot\mid s)).
\label{eq:ch2-mdpo-kl-gap}
\end{equation}
这里 KL 的左侧分布是当前 student，与 $\mathcal K_{k,n}^{S}=\operatorname{KL}( q_k\Vert p_{k,n})$ 的方向不同。

对 student 测度和 log-ratio 同时求导，有
\begin{align}
\nabla_\theta\mathcal L_{\mathrm{MDPO}}
&=\mathbb E_{d_k, p_{k,n}}\left[\left(\frac{U_{k,n}+1}\eta-f\right)\psi_{k,n}\right]\\
&=\mathbb E_{d_k, p_{k,n}}\left[\operatorname{sg}\!\left(\frac{U_{k,n}}\eta-f\right)\psi_{k,n}\right].
\label{eq:ch2-mdpo-gradient-bridge}
\end{align}
第二行使用 $\mathbb E_{ p_{k,n}}[\psi_{k,n}\mid s]=0$。其下降方向是 $-\nabla_\theta\mathcal L_{\mathrm{MDPO}}$，无需把 $\eta$ 吸收进 optimizer 步长。

\begin{proposition}[平方回归的投影含义]
\label{prop:ch2-regression-projection}
若 student 函数类不能表示式 \eqref{eq:ch2-logratio-equation} 的右侧，
PMD-Part 的精确 inner 解是镜像方程在函数类诱导的 $L_2(d_k q_k)$ 几何下的
投影；若训练样本的历史访问测度改为 $\nu$，投影也相应变为 $L_2(\nu q_k)$。
因此 inner loss 为零只在 target 可表示且训练测度覆盖目标测度时才意味着精确
mirror 解。
\end{proposition}

\begin{proof}
平方损失定义了可行 log-ratio 集合上的最小距离问题。对任意光滑可行曲线的一阶切向扰动 $v$，局部最优残差
$\mathscr E^\star$ 满足
\begin{equation*}
\mathbb E_{\nu, q_k}[\mathscr E^\star v]=0.
\end{equation*}
这是残差与可行切空间正交的必要条件。参数化可行集合通常非凸，该条件本身不保证全局最优。若 target 属于该集合，
投影残差为零；否则残差由函数类和测度共同决定。
\end{proof}
